\section{THẢO LUẬN KẾT QUẢ VÀ PHÂN TÍCH TƯƠNG QUAN}
\subsection{Chất lượng ảnh cao có bảo đảm nhận dạng tốt?}
Ở mức tối vừa của tập con, Zero-DCE đạt PSNR 20,08 dB, cao hơn diffusion 14,90 dB. Tuy nhiên, Macro-F1 của Zero-DCE chỉ đạt 69,30\%, thấp hơn cả Identity 75,35\%, trong khi diffusion đạt 82,71\%. Nếu chỉ chọn theo PSNR, kết luận sẽ không phù hợp với mục tiêu nhận dạng của đề bài.

Một ví dụ khác xuất hiện ở mức nặng: Retinexformer có PSNR 17,39 dB nhưng Macro-F1 52,57\%, còn diffusion có PSNR 14,88 dB và Macro-F1 63,72\%. Sự khác biệt cho thấy tiêu chí giữ cường độ gần ảnh gốc không bảo đảm giữ đúng những đặc trưng mà ResNet-18 sử dụng.

SSIM có xu hướng phù hợp hơn trong một số điều kiện. Ở mức nặng, diffusion dẫn đầu cả SSIM lẫn Macro-F1 trong tập con. Dù vậy, ở mức vừa CLAHE có SSIM 0,5568, cao hơn diffusion 0,5410, nhưng Macro-F1 thấp hơn. SSIM vì thế cũng chỉ là chỉ số bổ trợ, không phải hàm thay thế trực tiếp cho hiệu năng nhận dạng.

Về mặt cơ chế, các phép xử lý có thể tạo sai lệch độ sáng hoặc màu trên vùng lớn nhưng vẫn làm rõ nét chữ và viền. Ngược lại, một thuật toán có thể giảm sai số trung bình bằng cách làm mượt, đồng thời xóa mất chi tiết nhỏ. Đây là những cách giải thích hợp lý cho chênh lệch metric; để xác định nguyên nhân của từng mẫu cần quan sát ảnh và dự đoán ghép cặp, không thể suy ra chắc chắn chỉ từ bảng trung bình.

\subsection{Phương pháp phân tích tương quan}
Tương quan được tính riêng theo từng mức suy giảm. Mỗi phương pháp là một điểm gồm chất lượng ảnh trung bình và hiệu năng nhận dạng trung bình trên cùng 430 ảnh. Có tám điểm khi tính cả Identity. Clean bị loại khỏi phân tích vì PSNR của ảnh so với chính nó là vô hạn và vì đây không phải một phương án phục hồi ảnh tối.

Hệ số Pearson đo mức liên hệ tuyến tính giữa hai đại lượng:
\begin{equation}
r=\frac{\sum_i(x_i-\overline{x})(y_i-\overline{y})}
{\sqrt{\sum_i(x_i-\overline{x})^2}\sqrt{\sum_i(y_i-\overline{y})^2}}.
\end{equation}
Spearman được tính bằng Pearson trên thứ hạng, dùng thứ hạng trung bình nếu có giá trị bằng nhau. Hệ số này giúp xem thứ tự phương pháp theo chất lượng ảnh có phù hợp với thứ tự theo khả năng nhận dạng hay không, kể cả khi quan hệ không tuyến tính.

\begin{table}[htbp]
\centering
\caption{Tương quan chất lượng ảnh và nhận dạng của ResNet-18 trên tập con. Mỗi hàng có $n=8$ phương pháp, cùng mức suy giảm.}
\begin{tabular}{lllrr}
\toprule
Mức & Chất lượng & Nhận dạng & Pearson & Spearman \\
\midrule
\tableinput{tables/correlations}
\bottomrule
\end{tabular}
\label{tab:corr}
\end{table}

Các hệ số Pearson trong Bảng~\ref{tab:corr} đều dương, cho thấy chất lượng ảnh và hiệu năng nhận dạng có liên hệ cùng chiều khi xét toàn bộ tám phương pháp. Tuy nhiên, Spearman giữa PSNR và Macro-F1 giảm từ 0,667 ở mức nhẹ xuống 0,333 ở mức vừa và 0,286 ở mức nặng. Do đó, khi suy giảm mạnh hơn, xếp hạng theo PSNR trở nên ít phù hợp hơn với xếp hạng nhận dạng.

Với SSIM và Macro-F1, Spearman bằng 0,619 ở cả mức vừa và nặng, cao hơn tương ứng của PSNR. Quan sát này cho thấy thông tin cấu trúc hữu ích cho phân tích, nhưng không đủ để xác định một phương pháp tối ưu chỉ bằng SSIM. Các trường hợp đảo thứ hạng giữa CLAHE và diffusion vẫn tồn tại.

\fig{subset_scatter.png}{0.88}{Quan hệ giữa SSIM và Macro-F1 ở mức tối vừa, trích từ notebook ResNet-18. Mỗi điểm biểu diễn một phương pháp, không phải một ảnh độc lập.}{fig:scatter}

Số điểm trong phân tích rất nhỏ và các phương pháp sử dụng chung ảnh, nhiễu cũng như checkpoint. Vì vậy, các điểm không độc lập theo nghĩa thí nghiệm ngẫu nhiên. Một phương pháp có kết quả thấp ở cả hai trục, chẳng hạn MSR, có thể ảnh hưởng đáng kể đến Pearson. Báo cáo xem đây là phân tích khám phá, không suy ra quan hệ nhân quả và không dùng hệ số tương quan để bảo đảm hiệu quả trên ảnh đêm thực.

\subsection{Ablation: tách tác động của độ tối và nhiễu}
Để kiểm tra ảnh hưởng của từng thành phần suy giảm, notebook giữ $\gamma_d=2,2$ và khảo sát bốn biến thể: chỉ Gamma; Gamma kèm Gaussian với $\sigma_n=0,01$; Gamma kèm Poisson với $p=60$; Gamma kèm cả hai. Các bộ tăng cường giữ nguyên tham số đã chọn; ResNet-18 giữ nguyên checkpoint. Không chọn lại phương pháp riêng cho từng hàng ablation.

Ma trận phụ gồm $4\times3\times8\times3=288$ lượt trên 430 ảnh. Với trường hợp chỉ Gamma, thay noise seed không tạo ra nhiễu mới; các lượt tương ứng không được hiểu là những ảnh suy giảm độc lập. Bảng dưới trình bày trung bình Macro-F1 từ đầu ra notebook, làm tròn đến hai chữ số phần trăm.

\begin{table}[htbp]
\centering
\caption{Macro-F1 (\%) của ResNet-18 khi tách thành phần suy giảm ở mức Gamma 2,2.}
{\fontsize{11.5bp}{13.5bp}\selectfont
\begin{tabular}{lrrrr}
\toprule
Phương pháp & Chỉ Gamma & + Gaussian & + Poisson & + Cả hai \\
\midrule
\tableinput{tables/ablation}
\bottomrule
\end{tabular}}
\label{tab:ablation}
\end{table}

MSR đạt Macro-F1 98,20\% khi chỉ có Gamma, nhưng giảm xuống 73,33\% với Gaussian và 50,35\% với Poisson. Khi thêm cả hai thành phần nhiễu, kết quả còn 50,97\%, giảm 47,23 điểm phần trăm so với Gamma-only. Đây là bằng chứng rõ rằng phép xử lý hiệu quả với ảnh tối đơn thuần có thể không bền vững khi xuất hiện nhiễu.

Với Zero-DCE, Macro-F1 từ 97,48\% ở Gamma-only giảm còn 69,30\% khi có cả hai nhiễu. Diffusion giảm từ 95,04\% xuống 82,71\%, còn CLAHE giảm từ 93,30\% xuống 80,14\%. Trong cấu hình khảo sát, diffusion giữ hiệu năng tốt hơn khi nhiễu được đưa vào, dù không đứng đầu trường hợp chỉ Gamma.

\fig{ablation.pdf}{1}{Ảnh hưởng của từng thành phần nhiễu lên một số phương pháp đại diện. Biểu đồ dựng lại từ bảng ablation ResNet-18 trong notebook.}{fig:ablation}

Trong nhiều hàng, Poisson gây giảm F1 mạnh hơn Gaussian. Kết luận này chỉ đúng với mức $p=60$ và $\sigma_n=0,01$ đã chọn; hai loại nhiễu không được cân bằng về năng lượng nên không thể kết luận Poisson luôn nghiêm trọng hơn Gaussian. Ngoài ra, MSR ở trường hợp cả hai nhiễu cao hơn một chút so với chỉ Poisson, cho thấy tác động kết hợp không nhất thiết đơn điệu trên metric phân loại.

\subsection{Chi phí xử lý và khả năng triển khai}
\begin{table}[htbp]
\centering
\caption{Độ trễ trung bình của khối tăng cường ở mức vừa, trích từ 219 bản ghi notebook ResNet-18. Đơn vị ms/ảnh.}
\begin{tabular}{lr}
\toprule
Phương pháp & Thời gian \\
\midrule
\tableinput{tables/latency}
\bottomrule
\end{tabular}
\label{tab:latency}
\end{table}

Notebook ghi thiết bị GPU Tesla T4; các phép truyền thống trong mã nguồn chạy trên CPU, còn các mô hình sâu chạy trên thiết bị suy luận được cấu hình. Độ trễ được đo cho khối enhancer, không bao gồm toàn bộ quy trình phát hiện, cắt biển báo và phân loại. Mốc Identity vẫn có chi phí gọi hàm và quản lý dữ liệu nên thời gian đo không bằng 0.

CLAHE mất khoảng 0,400 ms/ảnh, còn diffusion khoảng 222,658 ms/ảnh ở mức vừa của notebook. Tỷ lệ chênh lệch xấp xỉ 557 lần, trong khi Macro-F1 khác khoảng 2,57 điểm phần trăm. So sánh này cho thấy cần xét ngân sách tính toán khi lựa chọn phương pháp, nhưng không được suy ra tốc độ toàn hệ thống chỉ từ thời gian enhancer.

Nếu yêu cầu ưu tiên chất lượng nhận dạng trong giao thức đã khảo sát, diffusion là lựa chọn phù hợp với quyết định validation. Nếu cần một phương án xử lý nhẹ, dễ triển khai, CLAHE là ứng viên thực tế để đánh giá tiếp. Việc thay đổi mục tiêu sang ràng buộc thời gian phải được quyết định và kiểm tra trên validation, không chọn lại sau khi xem test.

\subsection{Giới hạn của thực nghiệm}
Thứ nhất, ảnh thiếu sáng là ảnh tổng hợp trong miền RGB. Mô hình chưa mô tả đầy đủ phản ứng cảm biến, cân bằng trắng, nén ảnh, mưa, phản xạ gương, lóa đèn hoặc nhòe chuyển động. Khả năng khái quát đến camera giao thông ban đêm vì thế cần dữ liệu thực riêng.

Thứ hai, dữ liệu đầu vào đã được cắt biển báo. Kết quả chưa phản ánh lỗi phát hiện ROI trong ảnh toàn cảnh. Nếu detector bỏ sót hoặc cắt sai vị trí, bộ tăng cường crop không thể tự sửa được vấn đề đó.

Thứ ba, khảo sát chi tiết dùng 430 ảnh cân bằng lớp và ba training seed. Đây là quy mô phù hợp để khảo sát nhiều biến thể nhưng chưa đủ khẳng định ưu thế thống kê cho những khoảng cách rất nhỏ. Kết quả full-test được đối chiếu từ tài liệu dự án và cần bộ dự đoán gốc để kiểm định lại chênh lệch ghép cặp.

Thứ tư, lưới tham số còn hữu hạn và các bộ tăng cường sâu dùng trọng số đã học trên miền khác. Từ ``tối ưu'' trong báo cáo được hiểu là tốt nhất theo tiêu chí validation trong tập phương án đã khảo sát, không phải tối ưu tuyệt đối cho mọi dữ liệu và mọi thiết bị.
